\section{A single harmonic saddle throughout the moving range}
\label{us:sec:uniform}

For an integer $h\ge1$ define
\[
 e_h(n)=[u^h]\prod_{j=1}^n(1+u/j),\qquad
 T_{p,h}(a)=\sum_{n\ge0}\frac{(-1)^ne_h(n)}{(n+a)^p},
 \qquad p,a>0.
\]
The empty product is one, so $e_h(n)=0$ for $n<h$ and
$e_h(h)=1/h!$. In particular $e_1(n)=H_n$. The first nonzero
term of the alternating series motivates the normalization
\begin{equation}\label{us:eq:R}
 R_{p,h}(a)=(-1)^hh!(h+a)^pT_{p,h}(a).
\end{equation}
For $a=1/2$, the Gaussian harmonic family in the manuscript is
$S_{p,h}=2^{-p}T_{p,h}(1/2)$; its usual $S_p$ means $S_{p,1}$.

\subsection{The exact positive integral}

We briefly justify the starting point to keep the asymptotic argument
self-contained. The generating function is
\begin{equation}\label{us:eq:generating}
 \sum_{n\ge0}e_h(n)z^n
 =\frac{(-\log(1-z))^h}{h!(1-z)},\qquad |z|<1.
\end{equation}
It follows by taking the coefficient of $u^h$ in
$\sum_{n\ge0}\prod_{j=1}^n(1+u/j)z^n=(1-z)^{-1-u}$.
Ordinary convergence of $T_{p,h}$ for $p>0$ can also be checked
directly. The recurrence
$e_h(n)-e_h(n-1)=e_{h-1}(n-1)/n$ and
$e_h(n)=O_h((1+\log n)^h)$ show that
$e_h(n)/(n+a)^p\to0$ and its successive differences have an
absolutely summable majorant of order
$n^{-p-1}(1+\log n)^h$. Summation by parts against $(-1)^n$
therefore proves convergence. Absolute convergence holds for $p>1$.

Insert an Abel factor $0<\rho<1$ in the sum and use
$(n+a)^{-p}=\Gamma(p)^{-1}\int_0^\infty y^{p-1}e^{-(n+a)y}\,dy$.
Equation~\eqref{us:eq:generating} and dominated convergence as
$\rho\uparrow1$ give
\begin{equation}\label{us:eq:integral}
 R_{p,h}(a)=\frac{(h+a)^p}{\Gamma(p)}
 \int_0^\infty
  \frac{y^{p-1}e^{-ay}L(y)^h}{1+e^{-y}}\,dy,
 \qquad L(y)=\log(1+e^{-y}).
\end{equation}
The bound $L(y)\le\min(\log2,e^{-y})$ controls both endpoints.
This integral, including its positivity, is already established in
the manuscript \cite[\src{chapters/04-depth.tex},
\src{chapters/04-proportional-depth.tex}]{cont:ProveIt}.

Define
\[
 k_h(t)=\frac1{1+t}\left(\frac{\log(1+t)}t\right)^h,
 \quad k_h(0)=1.
\]
If $Y_p$ has the gamma law with shape $p$ and rate $h+a$, then
\begin{equation}\label{us:eq:gamma}
 R_{p,h}(a)=\mathbb E k_h(e^{-Y_p}).
\end{equation}
The kernel $k_h$ is strictly decreasing on $[0,1]$, with
$k_h(1)=(\log2)^h/2$ and $k_h(0)=1$. Consequently
\begin{equation}\label{us:eq:bounds}
 \frac{(\log2)^h}{2}<R_{p,h}(a)<1.
\end{equation}

\subsection{The moving saddle and the uniform expansion}

Put
\[
 w(y)=-\frac{d}{dy}\log L(y)
     =\frac1{(1+e^y)L(y)},\qquad
 \alpha=p/h,\qquad \epsilon=a/h.
\]
Writing $v=e^{-y}$ gives
\[
 \frac{d}{dv}\log\frac{v}{(1+v)\log(1+v)}
 =\frac{\log(1+v)-v}{v(1+v)\log(1+v)}<0.
\]
Thus $w$ increases strictly from $w_0=1/(2\log2)$ to $1$.
There is exactly one $x>0$ satisfying
\begin{equation}\label{us:eq:saddle}
 x(w(x)+\epsilon)=\alpha,
 \qquad
 \frac{p}{h+a}<x<\frac{p}{hw_0+a}.
\end{equation}
This is the saddle for the full exponent after absorbing the
factor $e^{-ay}$ into it. Keeping $a$ in the exponent is useful when
$a$ itself varies.

At this $x$ set
\begin{align}
 B&=1+\frac{x^2w'(x)}{\alpha},
 &A(z)&=\frac1{z(1+e^{-xz})},\label{us:eq:BA}\\
 \Psi(z)&=\log z+
  \frac{\log L(xz)-\log L(x)-\epsilon x(z-1)}\alpha.
 \label{us:eq:phase}
\end{align}
Then $\Psi(1)=\Psi'(1)=0$ and $\Psi''(1)=-B$.

\begin{theorem}[Uniform expansion with no ratio restriction]
\label{us:thm:uniform}
For each fixed integer $J\ge0$, as $p\to\infty$, uniformly over
all integers $h\ge1$ and all real $a>0$,
\begin{equation}\label{us:eq:uniform}
 R_{p,h}(a)=\mathcal S_{p,h}(a)
 \left\{\sum_{j=0}^J\frac{E_j(x,\epsilon)}{p^j}
       +O_J(p^{-J-1})\right\},
\end{equation}
where $E_0=1$ and
\begin{equation}\label{us:eq:prefactor}
 \mathcal S_{p,h}(a)=
 \frac{((h+a)x)^p e^{-ax}L(x)^h}
      {\Gamma(p)(1+e^{-x})}
 \sqrt{\frac{2\pi}{pB}}.
\end{equation}
Every coefficient $E_j$ is bounded on the entire parameter range.
In particular, with all derivatives evaluated at $z=1$,
\begin{equation}\label{us:eq:E1}
 E_1=\frac{A''}{2BA}
      +\frac{\Psi'''A'}{2B^2A}
      +\frac{\Psi''''}{8B^2}
      +\frac{5(\Psi''')^2}{24B^3}.
\end{equation}
The constants in the relative remainder depend on $J$ alone.
\end{theorem}

\begin{proof}
Substitute $y=xz$ in \eqref{us:eq:integral}. After removing its
value at the saddle, the remaining integral is
\begin{equation}\label{us:eq:reduced-integral}
 \int_0^\infty e^{p\Psi(z)}A(z)\,dz.
\end{equation}
We verify the uniform conditions for Laplace expansion directly.
The local expansion itself is the usual one
\cite[\S2.3(iii)]{cont:DLMFLaplace}; the global parameter control is
the point requiring proof here.

First, the concavity of $\log L$, equivalent to $w'>0$, gives its
tangent bound at $x$. Combining this bound with the saddle equation
yields the parameter-free majorant
\begin{equation}\label{us:eq:majorant}
 \boxed{\displaystyle \Psi(z)\le\log z-z+1\quad(z>0).}
\end{equation}
Also $A(z)\le z^{-1}$. For each fixed $0<d<1$, the contribution
of $|z-1|\ge d$ to \eqref{us:eq:reduced-integral} is bounded by
the corresponding tail of
$e^p\int_0^\infty z^{p-1}e^{-pz}\,dz$.
The gamma Chernoff bound, followed by Stirling's formula, bounds
that tail by $C_dp^{-1/2}e^{-c_dp}$ for positive constants
independent of $h$ and $a$.

Second, the saddle curvature is bounded above and below uniformly.
For $r_x=(1+e^x)^{-1}$, logarithmic differentiation gives
$w'=w(w-1+r_x)$, hence $0<w'\le r_x\le e^{-x}$.
Therefore
\begin{equation}\label{us:eq:curvature-bound}
 1\le B=1+\frac{xw'(x)}{w(x)+\epsilon}
 \le1+\frac{2\log2}{e},
 \qquad \frac12\le A(1)\le1.
\end{equation}
For $z$ in any fixed closed neighborhood of $1$ contained in
$(0,\infty)$, every fixed derivative of $A$ is uniformly bounded.
For $k\ge2$, the nonlogarithmic part of $\Psi^{(k)}(z)$ has
absolute value
\begin{equation}\label{us:eq:derivative-control}
 \frac{x^k|w^{(k-1)}(xz)|}{\alpha}
 =\frac{x^{k-1}|w^{(k-1)}(xz)|}{w(x)+\epsilon}.
\end{equation}
Near $x=0$, the analyticity of $w$ and $w\ge w_0$ bound this
expression. As $x\to\infty$, use
$w^{(j)}(u)=O_j(e^{-u})$ for $j\ge1$ and the boundedness of
$x^je^{-cx}$ for fixed $j$ and $c>0$. These observations also
prove the required bounds for the logistic amplitude. The linear
$\epsilon$ term has disappeared from all derivatives of order
at least two; increasing $\epsilon$ only increases the denominator
in \eqref{us:eq:derivative-control}. This explains why no upper
bound on $a$ is needed.

Finally,
$\Psi''(z)=-z^{-2}-x^2w'(xz)/\alpha\le-z^{-2}$.
On the fixed saddle neighborhood this supplies a common Gaussian
bound after $z=1+u/\sqrt p$. To any fixed order one may Taylor
expand the phase and amplitude there, retaining enough derivatives
to bound the remainder by $p^{-J-1}$ times a fixed integrable
polynomial times a Gaussian. One explicit justification is to first
restrict to $|u|\le p^\nu$, with a sufficiently small positive
$\nu$ depending on the desired order, expand to a higher finite
Taylor order, and use the common Gaussian bound on its complement.
All derivative bounds and all choices are independent of the
parameters. Extending each retained Gaussian monomial to the real
line incurs an exponentially small error. Odd monomials integrate
to zero. The even moments give
\[
 \int_0^\infty e^{p\Psi(z)}A(z)\,dz
 =A(1)\sqrt{\frac{2\pi}{pB}}
   \left\{\sum_{j=0}^J E_jp^{-j}+O_J(p^{-J-1})\right\}.
\]
Restoring the exact prefactor proves \eqref{us:eq:uniform}.
Taking the four possible even Gaussian contributions at order
$p^{-1}$ gives \eqref{us:eq:E1}.
\end{proof}

There is a finite prescription for every coefficient. Let
$\mathbb E_B$ denote expectation under a centered normal law of
variance $1/B$, and let $\tau$ be formal. Then
\begin{equation}\label{us:eq:coefficient-recipe}
 E_j=\frac1{A(1)}[\tau^{2j}]
 \mathbb E_B\!\left[
 A(1+\tau u)
 \exp\!\left(\sum_{k\ge3}
      \frac{\Psi^{(k)}(1)}{k!}u^k\tau^{k-2}\right)\right].
\end{equation}
Coefficient extraction precedes Gaussian integration. Only finitely
many derivatives and moments enter, so this is an exact algorithm
for any prescribed order, not an assertion that the complete
asymptotic series converges.

\subsection{Evaluation, overlaps, and limitations}

For the first correction, useful formulas are
\begin{align}
 \frac{A'(1)}{A(1)}&=-1+xr_x,&
 \frac{A''(1)}{A(1)}&=2-2xr_x+x^2(2r_x^2-r_x),\label{us:eq:amplitude-derivatives}\\
 \Psi'''(1)&=2-\frac{x^3w''(x)}\alpha,&
 \Psi''''(1)&=-6-\frac{x^4w'''(x)}\alpha.
 \label{us:eq:phase-derivatives}
\end{align}
As $x\to0$ or $x\to\infty$, $E_1\to1/12$ uniformly in
$\epsilon\ge0$. In the genuine first-pole limit, for example fixed
$h,a$ and $p\to\infty$, this correction cancels the first reciprocal
Stirling correction in the prefactor, consistently with $R\to1$.
The condition $x\to\infty$ alone is insufficient for that conclusion
when $h$ grows. For instance $p/h=\tfrac12\log h$, with fixed $a$,
still lies below the first-pole threshold and gives $R\to0$.

The theorem includes several formerly separate ranges:
$h=p^3$ gives $p/h\to0$; $h\asymp p$ gives a proportional
ratio; $p\asymp h\log h$ gives the first-pole transition; and fixed
$h$ gives an unbounded ratio. It also permits shifts such as
$a=p^2$. It does not cover a bounded outer exponent $p$. That
regime retains the separate endpoint expansion in the manuscript.
Nor does an $O_J$ theorem itself supply a certified numerical
constant at a specified small $p$.

For numerical work, compute $\log\mathcal S$ and standardize around
the saddle, rather than directly forming quantities which may be
exponentially small. The accompanying program uses
$u=\sqrt{pB}(z-1)$ in the exact integral. Table~\ref{us:tab:saddle}
reports selected checks; its last column is the absolute discrepancy
in the normalized bracket, $|R/\mathcal S-1-E_1/p|$.

\begin{table}[htbp]
\centering
\caption{Diagnostics for the uniformly scaled saddle formula.
Values come from 50-digit quadrature and are not interval bounds.}
\label{us:tab:saddle}
\begin{tabular}{@{}rrrrr@{}}
\toprule
$p$ & $h$ & $a$ & $|R/\mathcal S-1|$ & $|R/\mathcal S-1-E_1/p|$\\
\midrule
40 & $64000$ & $1/2$ & $2.08546\cdot10^{-3}$ & $2.12816\cdot10^{-6}$\\
160 & $4096000$ & $1/2$ & $5.20968\cdot10^{-4}$ & $1.34979\cdot10^{-7}$\\
40 & 40 & $1/2$ & $4.91662\cdot10^{-4}$ & $1.05027\cdot10^{-5}$\\
160 & 160 & $1/2$ & $1.15818\cdot10^{-4}$ & $6.80536\cdot10^{-7}$\\
160 & 1 & $1/2$ & $5.20968\cdot10^{-4}$ & $1.34979\cdot10^{-7}$\\
40 & 40 & 1600 & $2.08539\cdot10^{-3}$ & $2.12726\cdot10^{-6}$\\
\bottomrule
\end{tabular}
\end{table}

At $(p,h,a)=(160,4096000,1/2)$ the logarithm of the leading
prefactor is about $-1.501185\cdot10^6$. Relative asymptotics remain
informative even there. The finite test set is a check on signs,
normalizations, and the predicted order of error; the parameter-wide
statement follows from the proof, not from these samples.
